\section*{Part III: Choosing the Wedge}

\subsection*{28. A Beachhead Selection Is a Portfolio Decision Made Under Scarcity}

The founder does not choose the most interesting idea in isolation. The founder chooses \textbf{one use of scarce capital and attention instead of another.} That choice must consider what the company can actually do now, what it might learn, and the value of keeping other options open.

A segment can be attractive but poorly suited to the current founder. Another can be less profitable in the abstract but accessible through a credible relationship and easy to test with a weekend. The best first move may not maximize short-run revenue. It may maximize a combination of survival, evidence, useful capability, and the chance to repeat a successful transaction.

Do not replace these trade-offs with a fashionable acronym. Define what you need from the first business: perhaps paid demand within 3 months, positive cash contribution, work compatible with a full-time job, verified learner outcomes, and reusable academic assets. A different founder with investor capital, a large team, or substantial proprietary technology may rank opportunities differently.

\subsection*{29. The Entry Wedge Is a Mechanism, Not an Audience Description}

A company might have the right segment and still lack a way in. The wedge answers: \emph{Why would this buyer try us now despite the existence of acceptable alternatives and our lack of reputation?}

Consider several wedge mechanisms:

\begin{table*}[t]
\centering
\small
\begin{tabularx}{\textwidth}{l p{4.2cm} p{3.5cm} X}
\toprule
\textbf{Wedge} & \textbf{Why It Lowers the Barrier} & \textbf{Where It May Work} & \textbf{Central Danger} \\
\midrule
Urgent task & Fixes an immediate, consequential failure & Enterprise workflow, technical education & Urgency disappears or only rare buyers need it \\
Narrow specialization & Delivers much deeper relevance than a generalist & Consulting, vertical SaaS, education & Segment too small; expensive customization \\
Underserved service level & Adds reliable human attention or accountability & Clinics, education, B2B services & Labor costs exceed willingness to pay \\
Simplification & Removes costly features or complexity & Consumer tech, SMB software & Competitor can copy; remaining value too weak \\
Speed or convenience & Eliminates delay or coordination & Logistics, local services, software & Cost of speed and support underestimated \\
Risk reduction & Makes the initial commitment safer & Enterprise services, high-stakes purchases & Guarantee or pilot has unbounded liabilities \\
Channel access & Reaches buyers where competitors are absent & FMCG, partnerships, local services & Partner dependence, channel margin \\
Complement to incumbent & Works alongside current process without replacement & Developer tools, enterprise workflows & Never evolves into a valued standalone business \\
Specific use occasion & Fits a neglected moment of need & Consumer brands, travel, insurance & Occasion is infrequent or hard to predict \\
Manual concierge & Proves value before systems exist & Services, SaaS, marketplaces & Founder becomes permanent service bottleneck \\
\bottomrule
\end{tabularx}
\caption{Practical entry wedge mechanisms and their operational trade-offs.}
\label{tab:ch7-wedge-mechanisms}
\end{table*}

The same business can combine mechanisms, but the initial offer should usually emphasize 1 decisive reason to act rather than 10 vague promises.

\subsection*{30. The Wedge Must Cross 3 Adoption Thresholds}

A compelling first offer needs to cross 3 barriers:

\begin{itemize}
    \item \textbf{Attention:} The buyer recognizes that the product concerns an actual, immediate problem.
    \item \textbf{Trust:} The buyer believes the supplier might solve the problem without unacceptable downside.
    \item \textbf{Action:} The buyer can commit time, money, approvals, or a safe pilot under current constraints.
\end{itemize}

A technically brilliant feature may fail because buyers cannot explain it to their manager. An exceptionally cheap course may fail because employed learners cannot commit 6 evenings per week. A boutique consultancy may present credible expertise but demand a giant, open-ended contract from a cautious buyer.

Design the wedge around the largest adoption barrier that the founder can realistically change. Sometimes that means an inexpensive diagnostic, sometimes a precise guarantee, sometimes a partner introduction, sometimes a working integration, and sometimes a new product altogether.

\subsection*{31. Choose the Smallest Complete Job, Not the Smallest Fragment}

A common misunderstanding of starting small is producing a toy version that delivers none of the promised value. A restaurant does not test demand by serving photographs of food. An institute does not test transformation by sending a syllabus. An AI startup does not test workflow value through a polished chat interface that never completes a useful task.

\textbf{The ideal wedge completes 1 meaningful job end to end.} A 2-week applied studio might produce a tested evaluation harness and a defended technical design. A consultant's diagnostic might identify specific forecast failure causes and recommend an actionable decision. A software tool might reconcile 1 report with an audit trail.

Completeness creates evidence. Fragmentation creates opinions. If the customer experiences no meaningful result until the enterprise has purchased all 10 planned modules, the wedge may be too large to test safely.

\subsection*{32. Enter Through the Ignored Problem, Not Necessarily the Ignored Customer}

An overlooked segment can be visible, but the underserved \emph{job} may be more precise. Banks already buy fraud prevention. Factories already buy quality systems. Software teams already buy development tools. Yet a novel regulatory obligation, technical failure, or implementation bottleneck can create a specific new job that existing vendors address poorly.

Thus 2 paths exist: \textbf{new customers for an existing job} or \textbf{new jobs for an existing customer group}. The first may require building access; the second may borrow existing relationships. A company already serving colleges may find it easier to add applied labs than to sell to unrelated enterprises. A new founder with access to professional developers may find a focused engineering review easier than a national beginner-skilling program.

The choice should be based on observable transfer: relationships, trust, technical assets, payment mechanisms, or operations. ``Adjacent'' merely because both things use AI is not enough.

\subsection*{33. The Anti-Segment Is a Useful Strategic Tool}

A good beachhead definition should say whom you will \emph{not} serve, even if some excluded buyers are willing to pay. This prevents the first few sales from silently becoming a portfolio of unrelated obligations.

An early advanced engineering studio might decline complete beginners, people seeking placement guarantees, participants unable to devote sufficient practice time, and firms requiring access to confidential production code beyond the provider's security capability. The exclusions are not moral judgments. They protect the promised experience, economics, and risk boundaries.

Likewise, a new AI SaaS provider may refuse custom integrations that would consume 3 months for 1 customer. A specialist consulting firm may decline broad organizational transformation when its evidence and staff support only narrow diagnosis and implementation.

Write an \textbf{anti-customer profile} next to the target profile. Review it when tempting revenue appears.

\subsection*{34. Treat Customer Access as an Asset That Must Be Earned or Rented}

A group can have a perfect problem and still be inaccessible. An executive team in a regulated bank might need exactly what the founder can deliver, yet require long procurement, insurance, references, security certification, and relationships that the new company lacks.

Evaluate access separately:

\begin{enumerate}
    \item Can you name or locate individual prospective buyers?
    \item Can you approach them lawfully, respectfully, and without a costly mass campaign?
    \item Will they speak with an unknown provider?
    \item Does the buyer have power to act or a plausible path to the approver?
    \item Can 1 successful customer create a credible route to the next?
\end{enumerate}

A first beachhead requiring a 6-month sales cycle and several million rupees in vendor onboarding may be unsuitable for a founder with 3 months of cash, even if the potential contract is large.

\subsection*{35. The Founding Relationship Is Both Advantage and Bias}

Warm relationships can help a startup enter a difficult market. A former colleague knows the founder's work; a dean trusts the founder's teaching competence; an engineering manager knows they can deliver a real technical review.

But a warm relationship can also make demand seem stronger than it is. Friends purchase to be supportive. Existing contacts forgive weaknesses. The first large contract may rely on personal history that cannot be duplicated by a hired salesperson.

Mark each early customer as \textbf{relationship-led, referred, or independently acquired.} Analyze separately what made the purchase happen. A real relationship is an asset, but it is not evidence that strangers in the same segment will behave identically.
